\section{问题四：材料收缩模型与达标时间}
\subsection{质量基准和材料坐标}
使用附件2重构 $R(t)$，令 $\xi=r/R(t)$。均匀径向收缩假设给出材料速度 $v_r=\dot R\xi$，网格速度同为 $w_r=\dot R\xi$。干物质与水分守恒写为
\begin{align}
\partial_t\rho_d+\frac1r\partial_r(r\rho_dv_r)&=0,\label{eq:drymass}\\
\rho_d(\partial_tC+v_r\partial_rC)&=\frac1r\partial_r(r\rho_dD\partial_rC).\label{eq:wetmass}
\end{align}
初始 $\rho_d$ 均匀时，$\rho_d(t)R^2(t)=\rho_{d0}R_0^2$，故 $\rho_d$ 在每个时刻对空间均匀，可从\cref{eq:wetmass}中约去。这样固定的材料控制体具有固定干物质质量，$C$ 仍是水质量与干物质质量之比。

设 $\psi(\xi,t)=C(R(t)\xi,t)$，则
\begin{equation}
\left.\partial_tC\right|_r=\left.\partial_t\psi\right|_\xi-\frac{\dot R}R\xi\partial_\xi\psi,\qquad
v_r\partial_rC=\frac{\dot R}R\xi\partial_\xi\psi.\label{eq:chainmoving}
\end{equation}
两项相消，得到
\begin{align}
\partial_tC\big|_\xi&=\frac1{R^2\xi}\partial_\xi(\xi D\partial_\xi C),\label{eq:moving}\\
\rho(C)c_p(C)\partial_tT\big|_\xi&=\frac1{R^2\xi}\partial_\xi(\xi k\partial_\xi T).\label{eq:movingheat}
\end{align}
这里不能再单独加入 $-(\dot R/R)\xi\partial_\xi C$，否则相当于把网格运动当作材料与网格之间的相对运动。温度方程仍采用题给经验储热系数，并未由\cref{eq:drymass}重新解释为严格混合物能量守恒。

\subsection{边界、离散与质量核验}
表面条件变为
\begin{equation}
-\frac kR\partial_\xi T=h_T(T_s-T_e),\qquad
-\frac DR\partial_\xi C=h_m(C_s-C_e).\label{eq:movingbc}
\end{equation}
内部扩散尺度为 $R^{-2}$，表面项尺度为 $R^{-1}$。采用固定材料节点、Kirchhoff 界面通量及全隐式迭代，不需迎风项，也不使用 Courant 数作为精度证据。
记单位总干质量归一化的控制体权重为 $\omega_i=V_i/\sum_jV_j$，则每步检查
\begin{equation}
\varepsilon_n=\sum_i\omega_i(C_i^{n+1}-C_i^n)+\Delta t\frac{2h_m}{R^{n+1}}(C_s^{n+1}-C_e^{n+1}).\label{eq:balance}
\end{equation}
累计残差除以初始平均含水率，形成可比较的守恒指标。该检查与材料坐标的干物质质量权重一致。
\begin{figure}[!htbp]\centering\includegraphics[width=0.85\textwidth]{moving_boundary_solver.pdf}
\caption{材料坐标收缩求解与验证流程}\label{fig:moving}\end{figure}

\subsection{结果与同物性固定半径对照}
计算得到
\begin{equation}\widehat t_d=\QfourSeconds{}~\mathrm s=\QfourHours{}~\mathrm h,\qquad R(\widehat t_d)=\QfourRadius{}~\mathrm{cm}.\label{eq:q4time}\end{equation}
\begin{table}[!htbp]\centering
\caption{问题4每6 h及达标时含水率（kg/kg）}
\begin{tabular}{lrrrr}\toprule
$t$/h & 0 cm & 0.5 cm & 1.0 cm & 动态表面\\\midrule
6 & 1.7165 & 1.5350 & 1.0213 & 0.4216\\
12 & 0.7342 & 0.6517 & 0.4069 & 0.1666\\
18 & 0.4066 & 0.3668 & 0.2388 & 0.0889\\
24 & 0.2838 & 0.2599 & 0.1783 & 0.0671\\
30 & 0.2253 & 0.2085 & 0.1488 & 0.0593\\
36 & 0.1921 & 0.1790 & 0.1313 & 0.0558\\
42 & 0.1706 & 0.1598 & 0.1196 & 0.0539\\
48 & 0.1556 & 0.1463 & 0.1113 & 0.0529\\
50.8275 & 0.1500 & 0.1413 & 0.1081 & 0.0525\\
\bottomrule\end{tabular}\end{table}

表中固定物理位置通过 $\xi=r_j/R(t)$ 插值；当 $r_j>R(t)$ 时该位置位于材料外，留空；末列始终为 $\xi=1$ 的实时表面。result4.xlsx保留 $0,0.1,\ldots,1.9$~cm及动态表面，逐分钟输出至\QfourOutputEnd{}~s。

采用相同附录4物性、相同环境及相同数值设置，仅保持 $R=2$~cm的对照时间为\QfourFixedHours{}~h。收缩情景相对于该对照缩短\QfourReduction{}\%。这一比例是本模型及给定收缩曲线下的条件性比较，不代表对其他材料或收缩机理的普遍规律。
\begin{figure}[!htbp]\centering\includegraphics[width=0.94\textwidth]{q4_shrinkage_comparison.pdf}
\caption{同物性固定半径与材料收缩情景的比较}\label{fig:q4}\end{figure}
扩散距离缩短解释了主要加速方向。由于第三问采用另一组物性，不能将其与第四问的时间差全归因于收缩。本文也不再把纯坐标项的增删解释为独立的物理收缩机制。
